\documentclass[10pt,a4paper]{article}
\usepackage{amsmath,amssymb,amsthm,graphicx}
\usepackage{hyperref}
\usepackage{latexsym}
\usepackage{amssymb}
\usepackage{amsthm}
\usepackage{ulem}
\usepackage{enumitem}
\setlist{itemsep=0.4pt,topsep=0pt}
\parindent 0pt
\parindent 0pt
\parskip 12pt
\newcommand{\HCF}{\mathrm{hcf}}
\newcommand{\ZZ}{\mathbb{Z}}
\newcommand{\rk}{\mathrm{rk}}
\DeclareMathOperator{\len}{length}
\DeclareMathOperator{\bl}{b}
\DeclareMathOperator{\cl}{cl}
\DeclareMathOperator{\di}{d}
\DeclareMathOperator{\Cl}{Cl}
\DeclareMathOperator{\acl}{acl}
\DeclareMathOperator{\tp}{tp}
\DeclareMathOperator{\Th}{Th}
\newcommand{\Cor}{\mathcal{C}_\mathrm{or}}
\newcommand{\Cun}{\mathcal{C}_\mathrm{un}}
\newcommand{\mons}{\mathbb{M}}
\newcommand{\strong}{\mathrm{strong}}
\newcommand{\axioms}{\Delta_{\mathrm{FSTS}}}
\newtheoremstyle{note}% hnamei
{8pt}% hSpace abovei
{14pt}% hSpace belowi
{\normalfont}% hBody fonti
{}% hIndent amounti1
{\bfseries}% hTheorem head fonti
{ }% hPunctuation after theorem headi
{.5em}% hSpace after theorem headi2
{}% hTheorem head spec (can be left empty, meaning ‘normal’)i
\theoremstyle{note}
\newtheorem{thm}{Theorem}[section]
%\newtheorem{theorem}{Theorem}
\newtheorem{lemma}[thm]{Lemma} 
\newtheorem{df}[thm]{Definition}
\newtheorem{rmk}[thm]{Remark} 
\newtheorem{prop}[thm]{Proposition}
\newtheorem{fact}[thm]{Fact}
\newtheorem{cor}[thm]{Corollary}
\newtheorem{conj}[thm]{Conjecture}
\newtheorem{ex}[thm]{Example}
\newtheorem{qu}[thm]{Question}
\newtheorem*{notation}{Notation} 
\title{Every unconfined Steiner triple system carries an orientation}
\author{Silvia Barbina, Enrique Casanovas}
\date{arxiv:2411.13723v3; Creative Commons Attribution 4.0 International}
\begin{document}
\maketitle
\noindent\textit{Editorial note.} The counted statement is the author's environment at source0.tex lines 445-446 and its proof environment opens at line 447; the pairing is by adjacency and the retained proof range is the whole environment.\par\medskip
\noindent\textit{Editorial note.} The preamble is the author's own package lines, macro block and theorem-environment declarations (source lines 4-65); the author's optioned amsmath line and the fleqn-dependent layout line are not retained.\par\medskip
\noindent\textit{Editorial note.} The source carries no label-number dump and no counter values are imposed, so the printed numbering is the wrapper's default; identity is the author label, or the pinned source position where the author gave none.\par\medskip
\noindent\textit{Editorial note.} Metadata provenance: the source record of the pinned copy and the licence evidence file are both present; the licence is CC BY 4.0 with status supported, and the card's source\_version matches the value recorded in that evidence file.\par\medskip
\noindent\textit{Editorial scope.} One original result of the pinned author manuscript, counted once: the unconfined Steiner triple systems are oriented (unlabelled proposition at source0.tex lines 445-446).\par\medskip
\noindent\textit{Allowed standard prerequisites (editorial).}\begin{itemize}
\item Unconfined system. A Steiner triple system is unconfined when it carries an HF-ordering with no maximum.. It is the hypothesis of the inclusion.
\item Orientation. An orientation of A selects one member of each block as its representative, with the closure condition on the resulting relation.. It is the structure that has to be produced.
\item Inessential elements. An element belongs to at most one block of A when it is inessential, which makes it convenient to remove in the induction.. It is the case split of the first induction.
\end{itemize}\medskip
\begin{prop} $\Cun\subseteq \Cor \, .$
\end{prop}

\begin{proof}  We check by induction on $|A|$ that if $A$ is unconfined, then it has an orientation.  We may assume $A\neq \emptyset$. Let $a\in A$ belong to at most one block of $A$. By induction hypothesis, $A\smallsetminus \{a\}$ has an orientation $c$. If $a$ belongs to the block $e$, we make $a$ the apex of $e$. This gives an orientation of $A$ extending $c$.

Now we prove by induction on  $|A|$ that if $A$ is unconfined and  $X\subseteq A$, and $|X|\leq 2$, then $X\leq^* A$.  It is enough to check that if $X\subseteq A$, then $\delta(X)\geq \min\{ |X|,2\}$.  We may assume $A\neq\emptyset$. Choose $a\in A$ belonging to at most one block of $A$ and let $A^\prime= A\smallsetminus\{a\}$. We can apply the induction hypothesis to $A^\prime$. If $a\not\in X$, then $X\subseteq A^\prime$ and the induction hypothesis gives the result. If $a\in X$, let  $X^\prime =X\smallsetminus\{a\}$. Then  $\delta(X^\prime)=\delta(X)$ or  $\delta(X)=\delta(X^\prime)+1$.  If $|X|\leq 2$, then $\delta(X)=|X|$ and clearly $\delta(X)\geq \min\{|X|,2\}$.  If $|X|>2$, then %$\min\{|X|,2\}= 2$  and $\delta(X)\geq \delta(X^\prime)+1\geq 2$.
$|X^\prime| \geq 2$, and by induction hypothesis $\delta(X^\prime) \geq \min \{|X^\prime|,2 \}=2 $, hence $\delta(X) \geq  \delta(X^\prime) \geq 2= \min \{|X|, 2 \} $.
\end{proof}

\end{document}
